\documentclass[12pt,a4paper]{article}

% ---- Кириллица / шрифты ----
\usepackage{fontspec}
\setmainfont{Times New Roman}[
  BoldFont = Times New Roman Bold,
  ItalicFont = Times New Roman Italic,
  BoldItalicFont = Times New Roman Bold Italic
]
\setmonofont{Menlo}[Script=Default]
\usepackage{polyglossia}
\setdefaultlanguage{russian}
\setotherlanguage{english}
\newfontfamily\cyrillicfonttt{Menlo}[Script=Default]

% ---- Геометрия ----
\usepackage[left=25mm, right=15mm, top=20mm, bottom=20mm]{geometry}

% ---- Таблицы, списки, цвета ----
\usepackage{booktabs}
\usepackage{longtable}
\usepackage{tabularx}
\usepackage{array}
\usepackage{enumitem}
\usepackage{xcolor}
\usepackage{colortbl}

% ---- Ссылки, оглавление ----
\usepackage{hyperref}
\hypersetup{
  colorlinks=true,
  linkcolor=blue!70!black,
  urlcolor=blue!70!black,
  citecolor=blue!70!black,
  pdfauthor={KTZ Hackathon Team},
  pdftitle={ТЗ: Цифровой двойник локомотива v2.0}
}
\usepackage{titlesec}

% ---- Математика ----
\usepackage{amsmath}
\usepackage{amssymb}

% ---- Код ----
\usepackage{listings}
\lstset{
  basicstyle=\small\ttfamily,
  breaklines=true,
  frame=single,
  backgroundcolor=\color{gray!5},
  rulecolor=\color{gray!40},
  xleftmargin=2em,
  framexleftmargin=1.5em,
  numbers=none,
  tabsize=2
}

% ---- Колонтитулы ----
\usepackage{fancyhdr}
\pagestyle{fancy}
\fancyhf{}
\fancyhead[L]{\small\textit{ТЗ: Цифровой двойник локомотива v2.0}}
\fancyhead[R]{\small\textit{Хакатон КТЖ 2026}}
\fancyfoot[C]{\thepage}
\renewcommand{\headrulewidth}{0.4pt}

% ---- Цвета для статусов ----
\definecolor{statusgreen}{HTML}{22C55E}
\definecolor{statusyellow}{HTML}{EAB308}
\definecolor{statusred}{HTML}{EF4444}
\definecolor{sectionblue}{HTML}{1E40AF}

% ---- Форматирование секций ----
\titleformat{\section}{\Large\bfseries\color{sectionblue}}{\thesection}{1em}{}
\titleformat{\subsection}{\large\bfseries}{\thesubsection}{1em}{}
\titleformat{\subsubsection}{\normalsize\bfseries}{\thesubsubsection}{1em}{}

% ---- Утилиты ----
\newcommand{\code}[1]{\texttt{#1}}
\newcommand{\warn}[1]{\textcolor{statusred}{\textbf{#1}}}
\newcolumntype{L}[1]{>{\raggedright\arraybackslash}p{#1}}
\newcolumntype{C}[1]{>{\centering\arraybackslash}p{#1}}

\begin{document}

% =========================================================================
%  ТИТУЛЬНАЯ СТРАНИЦА
% =========================================================================
\begin{titlepage}
\centering
\vspace*{3cm}
{\Huge\bfseries ТЕХНИЧЕСКОЕ ЗАДАНИЕ\par}
\vspace{0.8cm}
{\LARGE Проект <<Цифровой двойник локомотива>>\par}
\vspace{0.4cm}
{\LARGE для хакатона КТЖ\par}
\vspace{2cm}
{\Large Версия 2.0\par}
\vspace{0.5cm}
{\large 04.04.2026\par}

\vfill

\begin{tabular}{ll}
\toprule
\textbf{Параметр} & \textbf{Значение} \\
\midrule
Заказчик & Хакатон КТЖ \\
Назначение & Рабочее ТЗ для команды разработки \\
Локомотивы & KZ8A (электровоз), ТЭ33А (тепловоз) \\
Стек & React + Next.js / FastAPI / TimescaleDB / Redis \\
Источник логики & \code{Locomotive\_Metrics\_Reference.md} \\
\bottomrule
\end{tabular}

\vspace{1.5cm}
{\small Документ фиксирует архитектурные, функциональные и нефункциональные требования.\\ Во всех местах, не определённых заказчиком жёстко, выбран наиболее подходящий инженерный вариант.}
\end{titlepage}

% =========================================================================
%  ОГЛАВЛЕНИЕ
% =========================================================================
\tableofcontents
\newpage

% =========================================================================
%  1. ОБЩАЯ ЦЕЛЬ ПРОЕКТА
% =========================================================================
\section{Общая цель проекта}

Разработать \textbf{production-like full-stack прототип} веб-приложения <<Цифровой двойник локомотива>>, которое в режиме, приближённом к реальному времени:

\begin{itemize}[nosep]
  \item принимает поток телеметрии от бортовых систем локомотива;
  \item рассчитывает интегральный \textbf{Health Index} (0--100) с объяснимостью (top-5 факторов);
  \item выявляет опасные отклонения и формирует алерты с рекомендациями;
  \item предоставляет машинисту, диспетчеру и администратору понятный интерфейс мониторинга, диагностики и анализа.
\end{itemize}

\textbf{Основной приоритет} --- безопасность локомотива, безопасность движения и снижение риска для окружающей среды.

При распределении усилий команда должна ориентироваться на \textbf{весовые критерии оценки} (см.~раздел~\ref{sec:criteria}): наибольший вес имеет работа в реальном времени (35\,\%) и дизайн интерфейса (30\,\%).


% =========================================================================
%  2. КОНТЕКСТ И ИСХОДНЫЕ ПРЕДПОСЫЛКИ
% =========================================================================
\section{Контекст и исходные предпосылки}

Реальные производственные данные КТЖ команде \textbf{недоступны}. Решение реализуется как реалистичный демонстрационный прототип, в котором бортовой источник заменяется \emph{lightweight simulator service}.

Демонстрационный пайплайн:

\begin{center}
\fbox{Simulator} $\;\longrightarrow\;$ \fbox{Redis Queue} $\;\longrightarrow\;$ \fbox{FastAPI Backend} $\;\longrightarrow\;$ \fbox{TimescaleDB} $\;\longrightarrow\;$ \fbox{WebSocket} $\;\longrightarrow\;$ \fbox{Next.js Frontend}
\end{center}

Все технические характеристики локомотивов и логика оценки параметров берутся из файла \code{Locomotive\_Metrics\_Reference.md}.

\subsection{Критерии соответствия безопасности}
\begin{itemize}[nosep]
  \item Не использовать реальные коммерческие данные без согласования.
  \item Обезличивание в примерах; mock-данные по умолчанию.
  \item Конфигурационные секреты --- \textbf{только через переменные окружения} (ENV).
\end{itemize}


% =========================================================================
%  3. КРИТЕРИИ ОЦЕНКИ ХАКАТОНА
% =========================================================================
\section{Критерии оценки хакатона}
\label{sec:criteria}

\begin{tabularx}{\textwidth}{C{0.6cm} L{3.8cm} C{1.2cm} X}
\toprule
\textbf{\#} & \textbf{Критерий} & \textbf{Вес} & \textbf{Что оценивается} \\
\midrule
1 & Дизайн интерфейса (UI/UX) & 30\,\% & Читабельность, приоритизация сигналов, информативные визуализации, accessibility, темизация. \\
\addlinespace
2 & Работа в реальном времени & 35\,\% & WebSocket/SSE, стабильность под нагрузкой, задержки, отказоустойчивость клиента и сервера. \\
\addlinespace
3 & Архитектура backend & 25\,\% & Чистая модульность, масштабируемость, хранение истории, прозрачность индекса здоровья. \\
\addlinespace
4 & Демонстрация и инженерная культура & 10\,\% & Чистота репозитория, инструкции запуска, логи/метрики, тестовый сценарий. \\
\bottomrule
\end{tabularx}


% =========================================================================
%  4. ПОДДЕРЖИВАЕМЫЕ ТИПЫ ЛОКОМОТИВОВ
% =========================================================================
\section{Поддерживаемые типы локомотивов}

Система обязана поддерживать \textbf{два типа} локомотивов:

\begin{tabularx}{\textwidth}{l l X}
\toprule
\textbf{Тип} & \textbf{Тяга} & \textbf{Особенности} \\
\midrule
KZ8A & Электровоз & Переменный ток 25 кВ, тяговые двигатели, рекуперативное торможение \\
ТЭ33А & Тепловоз & Дизель-электрическая, дизельный двигатель + генератор, расход топлива \\
\bottomrule
\end{tabularx}

\vspace{0.5em}
Для каждого типа существует \textbf{отдельный конфигурационный профиль}: список параметров, единицы измерения, безопасные/предупредительные/критические диапазоны, группировка по подсистемам, веса подсистем. Переключение типа локомотива \textbf{не требует изменения кода} --- используется конфигурационный подход.

Конкретные значения параметров и диапазонов~--- \textit{см.\,\code{Locomotive\_Metrics\_Reference.md}}.


% =========================================================================
%  5. СТЕК И ТЕХНИЧЕСКИЕ РЕШЕНИЯ
% =========================================================================
\section{Стек и принятые технические решения}

\begin{tabularx}{\textwidth}{l l X}
\toprule
\textbf{Слой} & \textbf{Технология} & \textbf{Комментарий} \\
\midrule
Frontend & React + Next.js & App Router, TypeScript, SSR для начального рендера \\
UI & Tailwind CSS + shadcn/ui & Быстрый стек с поддержкой темизации \\
Charts & ECharts & Высокопроизводительные интерактивные графики \\
Maps & Leaflet / MapLibre & Карта маршрута с маркерами и оверлеями \\
Backend & FastAPI (async) & Python 3.11+, асинхронные endpoints и WebSocket \\
Очередь & Redis / asyncio.Queue & Шина событий, буферизация потока телеметрии \\
Database & TimescaleDB & PostgreSQL + расширение для time-series, hypertables \\
Realtime & WebSocket (SSE fallback) & Поток обновлений backend $\to$ frontend \\
Infra & Docker Compose & Единый способ запуска всех сервисов \\
API Docs & OpenAPI / Swagger & Авто-генерация на \code{/docs} \\
Logging & structlog (JSON) & Структурированные логи, trace\_id \\
\bottomrule
\end{tabularx}


% =========================================================================
%  6. РОЛИ И АУТЕНТИФИКАЦИЯ
% =========================================================================
\section{Пользовательские роли и аутентификация}

\subsection{Роли}

\begin{tabularx}{\textwidth}{l X l}
\toprule
\textbf{Роль} & \textbf{Описание} & \textbf{Доступ к настройкам} \\
\midrule
Admin & Управление конфигурацией, порогами, сценарием симуляции, пользователями. Полный доступ ко всем экранам + панель администрирования. & Полный \\
\addlinespace
Dispatcher & Основной режим мониторинга: dashboard, alerts, trends, replay/history, export. Не может изменять пороги. & Только чтение \\
\addlinespace
Cabin mode & Упрощённый UI-режим для машиниста. Крупный Health Index, только ключевые показатели, высокий контраст. & Нет доступа \\
\bottomrule
\end{tabularx}

\subsection{Аутентификация и авторизация}

\begin{itemize}[nosep]
  \item \textbf{Механизм:} JWT-аутентификация (JSON Web Token).
  \item \textbf{Login screen:} форма с полями username/password.
  \item \textbf{Демо-пользователи:} предзаданные учётные записи в переменных окружения (\code{.env}):
    \begin{itemize}[nosep]
      \item \code{admin / admin123} --- роль Admin
      \item \code{dispatcher / disp123} --- роль Dispatcher
      \item \code{cabin / cabin123} --- роль Cabin
    \end{itemize}
  \item \textbf{Защита маршрутов:} frontend проверяет роль перед рендером; backend --- middleware на каждом endpoint.
  \item \textbf{Ограничение доступа:} Admin-only endpoints возвращают \code{403 Forbidden} для других ролей.
\end{itemize}


% =========================================================================
%  7. ФУНКЦИОНАЛЬНЫЕ ТРЕБОВАНИЯ
% =========================================================================
\section{Функциональные требования}

\subsection{Обязательные экраны MVP}

Login Page, Main Dashboard, Alerts Panel, Trends / Charts, Subsystem Cards, Карта маршрута, Replay / History, Report Export, Admin Settings.

\subsection{Login Page}

Центрированная карточка: логотип, поля username и password, кнопка <<Войти>>. При ошибке --- toast-уведомление. Автоматический редирект при наличии валидного токена.

\subsection{Main Dashboard}
\label{sec:dashboard}

\begin{itemize}[nosep]
  \item \textbf{Health Index виджет} --- крупный круговой индикатор (0--100), цветовая индикация статуса (Normal / Attention / Critical), текстовая подпись.
  \item \textbf{Top-5 факторов ухудшения} --- список параметров с наибольшим вкладом в снижение индекса. Для каждого: название, подсистема, текущее значение vs порог, полоса вклада.
  \item \textbf{Панель алертов} --- последние 10 алертов, иконка severity + время + сообщение, прокрутка.
  \item \textbf{Рекомендуемые действия} --- текстовый список, формируемый на основе активных алертов.
  \item \textbf{Панели ключевых узлов}:
    \begin{itemize}[nosep]
      \item \textbf{Скорость} --- текущая, предел, мини-график.
      \item \textbf{Топливо / Энергия} --- уровень/расход, тренд.
      \item \textbf{Давления / Температуры} --- ключевые показатели с цветовой индикацией.
      \item \textbf{Электрика} --- напряжение, ток, статус цепей.
    \end{itemize}
  \item \textbf{Селектор локомотива} --- выпадающий список (KZ8A / ТЭ33А).
  \item Автообновление через WebSocket; визуальный индикатор <<нет связи>> при потере соединения.
\end{itemize}

\subsection{Alerts Panel}
\label{sec:alerts}

\begin{itemize}[nosep]
  \item Таблица алертов: столбцы --- Время, Параметр, Severity (иконка + текст), Значение, Порог, Статус (new/acknowledged), Аннотация.
  \item Фильтры: чекбоксы severity, диапазон времени, селектор локомотива.
  \item Клик по строке --- разворачивает детали + поле для \textbf{отметки события} (аннотации): текстовое поле для ввода комментария к алерту/событию.
  \item Кнопка <<Acknowledge>> для подтверждения прочтения алерта.
  \item Пагинация.
\end{itemize}

\subsection{Карта маршрута / Схема пути}

\begin{itemize}[nosep]
  \item Упрощённая карта или схема участка пути (Leaflet/MapLibre).
  \item Маркер текущего положения локомотива.
  \item Наложение действующих скоростных ограничений (зоны на карте).
  \item Цветовая кодировка сегментов маршрута по состоянию Health Index.
  \item Overlay-панель: текущая скорость, координаты, ближайшая станция.
  \item Маршрут задаётся через GeoJSON файл в конфигурации.
\end{itemize}

\subsection{Trends / Charts}

\begin{itemize}[nosep]
  \item Интерактивные графики по ключевым метрикам (ECharts).
  \item Настраиваемое окно просмотра: 5, 10, 15 минут + произвольный диапазон.
  \item Tooltip с точными значениями, авто-скейлинг оси Y.
  \item Zoom (drag-select), pan, overlay нескольких метрик.
  \item Корректная работа с потоком данных от 1\,Гц и выше.
  \item Автообновление в реальном времени с возможностью паузы.
\end{itemize}

\subsection{Replay / History}

\begin{itemize}[nosep]
  \item Выбор временного диапазона (последние 5--15 минут).
  \item Timeline-скраббер (полоса перемотки) с кнопками Play / Pause.
  \item Управление скоростью воспроизведения: 1$\times$, 2$\times$, 5$\times$.
  \item \textbf{Синхронное обновление} всех компонентов: графики, Health Index, алерты, карта пути --- привязаны к позиции скраббера.
\end{itemize}

\subsection{Report Export}

\begin{itemize}[nosep]
  \item Экспорт мини-отчёта в формате \textbf{PDF} или \textbf{CSV}.
  \item Содержание отчёта: временной диапазон, сводка Health Index (min/avg/max), список сработавших алертов, средние значения ключевых метрик.
  \item Вызывается из экрана Replay или Dashboard.
\end{itemize}

\subsection{Admin Settings (только роль Admin)}

\begin{itemize}[nosep]
  \item Интерфейс с вкладками:
    \begin{itemize}[nosep]
      \item \textbf{Пороги} --- таблица по типу локомотива: параметр, safe\_low, safe\_high, warn\_low, warn\_high, crit\_low, crit\_high, weight. Инлайн-редактирование.
      \item \textbf{Симуляция} --- кнопки Start/Stop, выпадающий список сценариев, индикатор статуса.
      \item \textbf{Пользователи} --- простая CRUD-таблица (username, role, сброс пароля).
    \end{itemize}
\end{itemize}

\subsection{Адаптивность и доступность}

\begin{itemize}[nosep]
  \item Адаптивная вёрстка: корректное отображение на мониторах и ноутбуках различных разрешений.
  \item Обязательна \textbf{светлая и тёмная тема} с переключателем в header.
  \item Accessibility (WCAG~AA): достаточный контраст, читабельный размер шрифта, текстовые tooltips на всех иконках, навигация клавиатурой для критичных элементов.
\end{itemize}


% =========================================================================
%  8. АРХИТЕКТУРА СИСТЕМЫ
% =========================================================================
\section{Архитектура системы}

\subsection{Компоненты}

\subsubsection*{Simulator Service}
Отдельный Python-процесс (или контейнер), генерирующий телеметрию для KZ8A и ТЭ33А. Поддерживает нормальные и аварийные сценарии. Публикует JSON-сообщения в Redis Queue с частотой 1\,Гц на локомотив.

\subsubsection*{Шина событий (Redis / asyncio.Queue)}
Легковесная очередь, отвязывающая simulator от backend. Буферизирует входящий поток для защиты от всплесков нагрузки. Гарантирует порядок FIFO.

\subsubsection*{FastAPI Backend}
\begin{itemize}[nosep]
  \item Ingest worker: забирает сообщения из очереди, валидирует, нормализует, сглаживает.
  \item Health Index Engine: пересчитывает индекс на каждый тик.
  \item Alert Generator: проверяет пороги, создаёт алерты.
  \item REST API: исторические данные, конфигурация, экспорт.
  \item WebSocket Server: раздаёт live-обновления подключённым клиентам.
\end{itemize}

\subsubsection*{TimescaleDB}
PostgreSQL с расширением TimescaleDB. Хранение time-series телеметрии (hypertable), истории Health Index, алертов, конфигураций, пользователей. Retention policy 24--72 часа.

\subsubsection*{Next.js Frontend}
SPA на App Router + TypeScript. WebSocket-клиент с логикой переподключения. State management: React Context или Zustand. Рендеринг графиков через ECharts, карты через Leaflet.

\subsection{Архитектурная диаграмма}
\label{sec:arch-diagram}

В репозитории должна присутствовать архитектурная диаграмма (\code{docs/architecture.png} или \code{.svg}), отображающая:

\begin{itemize}[nosep]
  \item Все компоненты системы в рамках Docker Compose.
  \item Потоки данных: Simulator $\to$ Redis $\to$ FastAPI $\to$ TimescaleDB.
  \item Push-поток: FastAPI (WebSocket) $\to$ Browser.
  \item Pull-поток: Browser $\to$ FastAPI (REST) $\to$ TimescaleDB.
  \item Открытые порты (frontend: 3000, backend: 8000, DB: 5432).
\end{itemize}

\subsection{Логи, метрики, health-check}

\begin{itemize}[nosep]
  \item \textbf{Логи:} все backend-сервисы используют \code{structlog} (Python) для структурированных JSON-логов. Поля: \code{timestamp}, \code{level}, \code{component}, \code{message}, \code{trace\_id}.
  \item \textbf{Уровни логирования:} DEBUG (dev), INFO (prod по умолчанию), WARNING, ERROR.
  \item \textbf{Health-check endpoint:} \code{GET /health} --- возвращает статус сервиса, связность с БД, глубину очереди, uptime.
  \item \textbf{Метрики:} активные WebSocket-соединения, сообщений/сек (обработано), глубина очереди, время вычисления HI.
\end{itemize}


% =========================================================================
%  9. ОБРАБОТКА И ХРАНЕНИЕ ДАННЫХ
% =========================================================================
\section{Обработка и хранение данных}

\subsection{Data Pipeline}

Пошаговый поток обработки данных:

\begin{enumerate}[nosep]
  \item Simulator генерирует JSON-сообщение (\code{loco\_id}, \code{loco\_type}, \code{ts}, \code{parameters}).
  \item Сообщение помещается в Redis Queue (или asyncio.Queue).
  \item Ingest worker извлекает сообщение из очереди.
  \item \textbf{Валидация:} проверка обязательных полей, типов значений, формата timestamp.
  \item \textbf{Дедупликация:} отбрасываются сообщения с дублирующим \code{(loco\_id, ts)}.
  \item \textbf{Фильтрация выбросов:} значения за пределами физически возможного диапазона отвергаются (диапазоны настраиваются per parameter).
  \item \textbf{Обработка missing values:} forward-fill последним известным значением (максимум 5\,с, далее NaN).
  \item \textbf{Сглаживание:} EMA ($\alpha = 0.3$) или медианный фильтр (окно = 5). Выбирается per parameter в конфигурации.
  \item Запись raw + smoothed значений в \code{telemetry\_raw} (TimescaleDB).
  \item Триггер пересчёта Health Index $\to$ запись в \code{health\_index\_history} $\to$ push через WebSocket.
\end{enumerate}

\subsection{Схема TimescaleDB}
\label{sec:db-schema}

\subsubsection*{Таблица \code{telemetry\_raw} (hypertable)}

\begin{tabularx}{\textwidth}{l l X}
\toprule
\textbf{Поле} & \textbf{Тип} & \textbf{Описание} \\
\midrule
\code{ts} & \code{TIMESTAMPTZ NOT NULL} & Временная метка (partition key) \\
\code{loco\_id} & \code{VARCHAR(32) NOT NULL} & Идентификатор локомотива \\
\code{loco\_type} & \code{VARCHAR(16) NOT NULL} & Тип: KZ8A \textbar{} TE33A \\
\code{parameter} & \code{VARCHAR(64) NOT NULL} & Название параметра \\
\code{value\_raw} & \code{DOUBLE PRECISION} & Сырое значение \\
\code{value\_smooth} & \code{DOUBLE PRECISION} & Сглаженное значение (EMA / медиана) \\
\code{unit} & \code{VARCHAR(16)} & Единица измерения \\
\code{quality} & \code{SMALLINT DEFAULT 0} & 0 = ok, 1 = interpolated, 2 = outlier \\
\bottomrule
\end{tabularx}

\vspace{0.3em}
\noindent Chunk interval: 1 час. Retention policy: 72 часа (\code{drop\_chunks}).\\
Индекс: \code{(loco\_id, parameter, ts DESC)}.

\subsubsection*{Таблица \code{health\_index\_history} (hypertable)}

\begin{tabularx}{\textwidth}{l l X}
\toprule
\textbf{Поле} & \textbf{Тип} & \textbf{Описание} \\
\midrule
\code{ts} & \code{TIMESTAMPTZ NOT NULL} & Временная метка \\
\code{loco\_id} & \code{VARCHAR(32) NOT NULL} & Идентификатор локомотива \\
\code{score} & \code{DOUBLE PRECISION NOT NULL} & Значение HI (0--100) \\
\code{status} & \code{VARCHAR(16) NOT NULL} & Normal \textbar{} Attention \textbar{} Critical \\
\code{top\_factors} & \code{JSONB} & Массив top-5 факторов \\
\code{subsystem\_scores} & \code{JSONB} & Оценки подсистем \\
\bottomrule
\end{tabularx}

\vspace{0.3em}
\noindent Chunk interval: 1 час. Retention: 72 часа.

\subsubsection*{Таблица \code{alerts}}

\begin{tabularx}{\textwidth}{l l X}
\toprule
\textbf{Поле} & \textbf{Тип} & \textbf{Описание} \\
\midrule
\code{id} & \code{BIGSERIAL PK} & Автоинкремент ID \\
\code{ts} & \code{TIMESTAMPTZ NOT NULL} & Время срабатывания \\
\code{loco\_id} & \code{VARCHAR(32) NOT NULL} & Идентификатор локомотива \\
\code{parameter} & \code{VARCHAR(64) NOT NULL} & Параметр, вызвавший алерт \\
\code{severity} & \code{VARCHAR(16) NOT NULL} & info \textbar{} warning \textbar{} critical \\
\code{value} & \code{DOUBLE PRECISION} & Фактическое значение \\
\code{threshold} & \code{DOUBLE PRECISION} & Превышенный порог \\
\code{message} & \code{TEXT} & Текст алерта \\
\code{acknowledged} & \code{BOOLEAN DEFAULT FALSE} & Подтверждён ли \\
\code{annotation} & \code{TEXT} & Отметка события (аннотация) \\
\bottomrule
\end{tabularx}

\vspace{0.3em}
\noindent Индекс: \code{(loco\_id, ts DESC, severity)}.

\subsubsection*{Таблица \code{loco\_config}}

\begin{tabularx}{\textwidth}{l l X}
\toprule
\code{loco\_type} & \code{VARCHAR(16) PK} & Тип локомотива \\
\code{config} & \code{JSONB NOT NULL} & Параметры, диапазоны, веса, подсистемы \\
\code{updated\_at} & \code{TIMESTAMPTZ} & Время последнего обновления \\
\bottomrule
\end{tabularx}

\subsubsection*{Таблица \code{users}}

\begin{tabularx}{\textwidth}{l l X}
\toprule
\code{id} & \code{SERIAL PK} & ID пользователя \\
\code{username} & \code{VARCHAR(64) UNIQUE} & Логин \\
\code{password\_hash} & \code{VARCHAR(256)} & Хеш пароля (bcrypt) \\
\code{role} & \code{VARCHAR(16)} & admin \textbar{} dispatcher \textbar{} cabin \\
\bottomrule
\end{tabularx}

\subsection{Retention и агрегация}

\begin{itemize}[nosep]
  \item Retention сырых данных: \textbf{72 часа} (настраивается через ENV).
  \item Для replay --- прямые запросы к hypertable по временному диапазону.
  \item Опционально: continuous aggregate --- 1-минутные средние для графиков с большим окном.
\end{itemize}


% =========================================================================
%  10. HEALTH INDEX — МАТЕМАТИЧЕСКАЯ МОДЕЛЬ
% =========================================================================
\section{Индекс здоровья --- математическая модель}
\label{sec:health-index}

\subsection{Обзор}

Health Index (HI) --- единый числовой показатель $\in [0, 100]$, отражающий общее состояние локомотива. Расчёт \textbf{иерархический}:

\[
\text{Параметры} \;\xrightarrow{\text{нормализация}}\; \text{Подсистемы} \;\xrightarrow{\text{взвешивание}}\; \text{Глобальный HI}
\]

Индекс учитывает: текущие (сглаженные) значения, тренды, активные диагностические сигналы и штрафы за алерты.

\subsection{Функция нормализации}

Для каждого параметра $p$ с текущим сглаженным значением $v$ определяется нормализованная оценка $\mathrm{norm}(v, p) \in [0, 1]$:

\[
\mathrm{norm}(v, p) = \begin{cases}
  1.0, & \text{если } v \in [\text{safe\_low}_p,\; \text{safe\_high}_p] \\[6pt]
  \displaystyle 1 - \frac{|v - s_{\text{near}}|}{t_{\text{bound}} - s_{\text{near}}}, & \text{если } v \in \text{warning zone} \\[10pt]
  0.0, & \text{если } v \notin [\text{crit\_low}_p,\; \text{crit\_high}_p]
\end{cases}
\]

где $s_{\text{near}}$ --- ближайшая граница safe-зоны, $t_{\text{bound}}$ --- соответствующая граница critical-зоны.

Особые случаи: если $v = \mathrm{NaN}$ (пропуск данных $> 5$\,с), то $\mathrm{norm} = 0$ (Critical).

Конкретные значения \code{safe\_low}, \code{safe\_high}, \code{warn}, \code{crit} для каждого параметра --- \textit{см.\,\code{Locomotive\_Metrics\_Reference.md}}.

\subsection{Штрафы за алерты}
\label{sec:alert-penalty}

Активные алерты накладывают мультипликативный штраф на оценку затронутой подсистемы:

\[
\mathrm{penalty}(\text{sub}) = \max\!\Big(0.5,\;\; 1 - 0.1 \cdot N_{\text{crit}} - 0.05 \cdot N_{\text{warn}}\Big)
\]

где $N_{\text{crit}}$ --- количество активных critical-алертов в подсистеме, $N_{\text{warn}}$ --- warning-алертов. Минимальный множитель --- 0.5 (максимальный штраф --- 50\,\%).

\subsection{Оценка подсистемы}

Каждый тип локомотива определяет набор подсистем (например, для KZ8A: Тяга, Электрика, Охлаждение, Тормоза, Вспомогательные). Оценка подсистемы $\text{sub}$:

\[
S_{\text{sub}} = \mathrm{penalty}(\text{sub}) \cdot \frac{\displaystyle\sum_{p \in \text{sub}} w_p \cdot \mathrm{norm}(v_p, p)}{\displaystyle\sum_{p \in \text{sub}} w_p}
\]

где $w_p$ --- вес параметра $p$ внутри подсистемы (из конфигурации). Конкретные подсистемы и веса --- \textit{см.\,\code{Locomotive\_Metrics\_Reference.md}}.

\subsection{Глобальный Health Index}

\[
\mathrm{HI} = \left\lfloor \frac{\displaystyle\sum_{\text{sub}} W_{\text{sub}} \cdot S_{\text{sub}}}{\displaystyle\sum_{\text{sub}} W_{\text{sub}}} \cdot 100 \right\rceil
\]

где $W_{\text{sub}}$ --- вес подсистемы в глобальном индексе. Результат округляется до 1 знака после запятой и ограничивается $[0, 100]$.

\subsection{Категоризация}

\begin{tabularx}{\textwidth}{C{2.5cm} C{2.5cm} C{2cm} X}
\toprule
\textbf{Диапазон} & \textbf{Статус} & \textbf{Цвет} & \textbf{Описание} \\
\midrule
70 -- 100 & Normal & \textcolor{statusgreen}{\rule{1cm}{0.4cm}} & Все системы в безопасных диапазонах \\
40 -- 69 & Attention & \textcolor{statusyellow}{\rule{1cm}{0.4cm}} & Один или более параметров в warning-зоне \\
0 -- 39 & Critical & \textcolor{statusred}{\rule{1cm}{0.4cm}} & Требуется немедленное вмешательство \\
\bottomrule
\end{tabularx}

\subsection{Объяснимость --- Top-5 факторов}

После вычисления HI определяются 5 параметров с наибольшим негативным вкладом:

\[
\mathrm{contribution}(p) = W_{\text{sub}(p)} \cdot w_p \cdot \big(1 - \mathrm{norm}(v_p, p)\big)
\]

Сортировка по убыванию, берутся top-5. Каждая запись содержит: название параметра, подсистему, текущее значение, безопасный диапазон, величину вклада, рекомендуемое действие.

\subsection{Частота и производительность}

HI пересчитывается на каждый тик телеметрии (1\,Гц). Целевое время вычисления: $< 50$\,мс на один локомотив. Результаты кешируются в оперативной памяти и записываются в \code{health\_index\_history}.


% =========================================================================
%  11. СПЕЦИФИКАЦИЯ REST API
% =========================================================================
\section{Спецификация REST API}
\label{sec:api}

\subsection{Общие соглашения}

\begin{itemize}[nosep]
  \item Base URL: \code{/api/v1/}
  \item Content-Type: \code{application/json}
  \item Аутентификация: Bearer JWT в заголовке \code{Authorization}
  \item Ошибки: \code{\{"error": "code", "message": "...", "detail": \{...\}\}}
  \item Пагинация: \code{?limit=N\&offset=M}, заголовок ответа \code{X-Total-Count}
  \item Timestamps: ISO\,8601 с timezone (UTC)
\end{itemize}

\subsection{Authentication}

\begin{tabularx}{\textwidth}{l l X l}
\toprule
\textbf{Метод} & \textbf{Путь} & \textbf{Описание} & \textbf{Auth} \\
\midrule
POST & \code{/auth/login} & Аутентификация. Body: \code{\{username, password\}}. Ответ: \code{\{token, role, expires\_at\}} & --- \\
GET & \code{/auth/me} & Текущий пользователь: \code{\{id, username, role\}} & Required \\
\bottomrule
\end{tabularx}

\subsection{Telemetry}

\begin{tabularx}{\textwidth}{l l X l}
\toprule
\textbf{Метод} & \textbf{Путь} & \textbf{Описание} & \textbf{Auth} \\
\midrule
GET & \code{/telemetry/\{loco\_id\}/latest} & Последние значения всех параметров & dispatcher+ \\
GET & \code{/telemetry/\{loco\_id\}/history} & Исторические данные. Query: \code{parameter, from, to, interval} & dispatcher+ \\
POST & \code{/telemetry/ingest} & Ручная отправка телеметрии (тестирование) & admin \\
\bottomrule
\end{tabularx}

\subsection{Health Index}

\begin{tabularx}{\textwidth}{l l X l}
\toprule
\textbf{Метод} & \textbf{Путь} & \textbf{Описание} & \textbf{Auth} \\
\midrule
GET & \code{/health-index/\{loco\_id\}/current} & Текущий HI + top factors + subsystem scores & dispatcher+ \\
GET & \code{/health-index/\{loco\_id\}/history} & История HI. Query: \code{from, to} & dispatcher+ \\
\bottomrule
\end{tabularx}

\subsection{Alerts}

\begin{tabularx}{\textwidth}{l l X l}
\toprule
\textbf{Метод} & \textbf{Путь} & \textbf{Описание} & \textbf{Auth} \\
\midrule
GET & \code{/alerts/\{loco\_id\}} & Список алертов. Query: \code{severity, from, to, limit, offset} & dispatcher+ \\
PATCH & \code{/alerts/\{id\}/acknowledge} & Подтвердить прочтение алерта & dispatcher+ \\
PATCH & \code{/alerts/\{id\}/annotate} & Добавить отметку события. Body: \code{\{text\}} & dispatcher+ \\
\bottomrule
\end{tabularx}

\subsection{Configuration}

\begin{tabularx}{\textwidth}{l l X l}
\toprule
\textbf{Метод} & \textbf{Путь} & \textbf{Описание} & \textbf{Auth} \\
\midrule
GET & \code{/config/\{loco\_type\}} & Конфигурация локомотива (параметры, пороги, веса) & dispatcher+ \\
PUT & \code{/config/\{loco\_type\}} & Обновить конфигурацию & admin \\
GET & \code{/config/subsystems/\{loco\_type\}} & Список подсистем и их веса & dispatcher+ \\
\bottomrule
\end{tabularx}

\subsection{Simulator Control}

\begin{tabularx}{\textwidth}{l l X l}
\toprule
\textbf{Метод} & \textbf{Путь} & \textbf{Описание} & \textbf{Auth} \\
\midrule
POST & \code{/simulator/start} & Запустить симуляцию & admin \\
POST & \code{/simulator/stop} & Остановить симуляцию & admin \\
POST & \code{/simulator/scenario} & Сменить сценарий. Body: \code{\{scenario\}} & admin \\
GET & \code{/simulator/status} & Текущий статус симулятора & admin \\
\bottomrule
\end{tabularx}

\subsection{Export}

\begin{tabularx}{\textwidth}{l l X l}
\toprule
\textbf{Метод} & \textbf{Путь} & \textbf{Описание} & \textbf{Auth} \\
\midrule
GET & \code{/export/\{loco\_id\}} & Экспорт отчёта. Query: \code{format=pdf|csv, from, to} & dispatcher+ \\
\bottomrule
\end{tabularx}

\subsection{System}

\begin{tabularx}{\textwidth}{l l X l}
\toprule
\textbf{Метод} & \textbf{Путь} & \textbf{Описание} & \textbf{Auth} \\
\midrule
GET & \code{/health} & Health-check (DB, queue, uptime) & --- \\
GET & \code{/docs} & Swagger UI (авто-генерация) & --- \\
\bottomrule
\end{tabularx}

\subsection{WebSocket Protocol}

\textbf{Endpoint:} \code{ws://host/ws/telemetry/\{loco\_id\}}

\begin{enumerate}[nosep]
  \item Клиент подключается с JWT-токеном как query parameter (\code{?token=...}).
  \item Сервер аутентифицирует и начинает стриминг.
  \item Сервер отправляет JSON-фреймы с частотой 1\,Гц:
\end{enumerate}

\begin{lstlisting}
{
  "type": "telemetry_update",
  "data": {
    "ts": "2026-04-04T12:00:00.000Z",
    "loco_id": "KZ8A-001",
    "parameters": { "speed_kmh": 87.3, ... },
    "health_index": { "score": 82.5, "status": "Normal", ... },
    "new_alerts": [ ... ]
  }
}
\end{lstlisting}

\begin{enumerate}[nosep, start=4]
  \item Клиент может отправить: \code{\{"type": "ping"\}}, сервер ответит \code{\{"type": "pong"\}}.
  \item При отключении клиент реализует \textbf{exponential backoff}: 1\,с $\to$ 2\,с $\to$ 4\,с $\to$ 8\,с $\to$ max 30\,с.
  \item Визуальная индикация <<нет связи>> на UI в течение 3\,с после потери соединения.
  \item При реконнекте сервер сначала отправляет snapshot текущего состояния, затем возобновляет стриминг.
\end{enumerate}


% =========================================================================
%  12. ОПИСАНИЕ ЭКРАННЫХ ФОРМ (WIREFRAMES)
% =========================================================================
\section{Описание экранных форм}

\subsection{Login Screen}

Центрированная карточка на нейтральном фоне. Логотип проекта сверху, поля <<Username>> и <<Password>>, кнопка <<Войти>>. Toast-уведомление при ошибке. Автоматический редирект на Dashboard при наличии валидного токена. Минимальный дизайн, без sidebar.

\subsection{Main Dashboard}

\begin{itemize}[nosep]
  \item \textbf{Header} (full-width): логотип, селектор локомотива (dropdown), текущий пользователь + badge роли, переключатель темы, кнопка Logout.
  \item \textbf{Основная область}: Health Index hero-виджет (круговой gauge, число, текст статуса + цвет). Сетка панелей подсистем (Скорость, Топливо/Энергия, Давления/Температуры, Электрика). Каждая панель: мини-gauge + текущее значение + стрелка тренда ($\uparrow\downarrow\rightarrow$).
  \item \textbf{Боковая/дополнительная область}: список Top-5 факторов (параметр, bar вклада, значение vs порог). Лента алертов (scroll, последние 10). Рекомендации (текстовый список).
  \item \textbf{Нижняя полоса}: мини-график HI за последние 5\,мин, индикатор соединения.
\end{itemize}

\subsection{Alerts Panel}

Полноширинная таблица: столбцы --- Время, Параметр, Severity (иконка + текст), Значение, Порог, Статус, Аннотация. Над таблицей --- фильтры (severity checkboxes, date range picker, селектор локомотива). Клик по строке --- expand с деталями + текстовое поле аннотации. Пагинация внизу.

\subsection{Route Map}

Карта занимает основную часть экрана. Floating overlay: текущая скорость, координаты, ближайшая станция. Линия маршрута с цветовой кодировкой (зелёный/жёлтый/красный по HI). Маркер-иконка локомотива. Зоны ограничения скорости --- dashed overlay.

\subsection{Trends / Charts}

Панель выбора диапазона: кнопки-пресеты (5\,мин, 10\,мин, 15\,мин) + произвольный picker. Сетка графиков (настраиваемые метрики, количество и расположение --- на усмотрение дизайнера). Каждый график: line chart, tooltip, авто-скейлинг Y, zoom через drag-select. Кнопка <<Добавить метрику>>. Переключатель live/pause.

\subsection{Replay / History}

Timeline-скраббер (full-width) с кнопками Play / Pause / Speed (1$\times$, 2$\times$, 5$\times$). Под скраббером --- тот же layout, что и Dashboard, но данные управляются позицией курсора. Графики, карта, алерты, HI --- всё синхронизировано со скраббером.

\subsection{Admin Settings}

Tabs-интерфейс:
\begin{itemize}[nosep]
  \item \textbf{Пороги:} таблица per locomotive type. Строки = параметры, столбцы = safe\_low, safe\_high, warn\_low, warn\_high, crit\_low, crit\_high, weight. Inline-editing с валидацией.
  \item \textbf{Симуляция:} кнопки Start/Stop, dropdown сценариев (normal, overheat, fault, highload, connection\_loss), индикатор текущего статуса.
  \item \textbf{Пользователи:} простая CRUD-таблица (username, role, reset password).
\end{itemize}


% =========================================================================
%  13. СИМУЛЯТОР ТЕЛЕМЕТРИИ
% =========================================================================
\section{Симулятор телеметрии}

\subsection{Normal Operation Mode}

Random walk вокруг номинальных значений для каждого параметра, ограниченный safe-диапазонами. Гауссов шум с настраиваемой $\sigma$. Частота эмиссии --- 1\,Гц на локомотив.

\subsection{Аварийные сценарии (минимум 3)}

\begin{tabularx}{\textwidth}{C{0.6cm} l C{1.5cm} X}
\toprule
\textbf{\#} & \textbf{Сценарий} & \textbf{Длительность} & \textbf{Описание} \\
\midrule
1 & Постепенный перегрев & 2\,мин & Температура тягового двигателя дрейфует вверх, проходя warning $\to$ critical зоны. \\
2 & Внезапный электрический сбой & мгновенно & Напряжение падает до 0 за 1 тик. Множественные алерты. \\
3 & Highload burst & 30\,сек & Частота сообщений увеличивается до $\times$10. Тест backpressure. \\
\bottomrule
\end{tabularx}

\subsection{Connection Loss}

Симулятор прекращает отправку на 10--15 секунд, затем возобновляет. Frontend должен отобразить <<нет связи>> и корректно восстановиться.

\subsection{Конфигурация}

Сценарии описываются в YAML/JSON файлах. Каждый сценарий определяет: затронутые параметры, тип изменения (drift/step), скорость, длительность, целевые значения.


% =========================================================================
%  14. НЕФУНКЦИОНАЛЬНЫЕ ТРЕБОВАНИЯ
% =========================================================================
\section{Нефункциональные требования}

\subsection{Производительность / Latency}

\begin{itemize}[nosep]
  \item End-to-end latency (от эмиссии события до видимого обновления UI): $< 500$\,мс.
  \item Вычисление Health Index: $< 50$\,мс на локомотив на тик.
  \item Время ответа API (исторические запросы, 1-часовой диапазон): $< 200$\,мс.
  \item Frontend-рендеринг: 60\,fps в нормальном режиме.
\end{itemize}

\subsection{Highload / Burst}

\begin{itemize}[nosep]
  \item Система выдерживает $\times$10 burst (10\,Гц вместо 1\,Гц) без аварийной остановки.
  \item Очередь буферизирует всплески; backend обрабатывает на устойчивой скорости.
  \item Frontend может снизить частоту рендеринга до 2--4\,fps при burst, но \textbf{не должен зависать}.
  \item Graceful degradation: при отставании пропускаются промежуточные фреймы, всегда отображается последнее состояние.
\end{itemize}

\subsection{Отказоустойчивость}

\begin{itemize}[nosep]
  \item WebSocket reconnect с exponential backoff (1, 2, 4, 8, max 30\,с).
  \item Визуальная индикация <<нет связи>> на frontend в течение 3\,с после disconnect.
  \item Backend: при временной недоступности БД --- отдача данных из in-memory кеша до 30\,с.
  \item Нет потери данных в очереди при перезапуске backend (Redis persistence или graceful shutdown asyncio.Queue).
\end{itemize}

\subsection{Безопасность}

\begin{itemize}[nosep]
  \item JWT-аутентификация для всех API endpoints (кроме \code{/health}, \code{/docs}).
  \item Role-based access control: Admin --- полный доступ, Dispatcher --- read-only, Cabin --- read-only subset.
  \item Секреты (пароль БД, JWT secret) --- \textbf{только в переменных окружения}, никогда в коде.
  \item Все данные --- mock/обезличенные.
\end{itemize}

\subsection{Документация API}

\begin{itemize}[nosep]
  \item Авто-генерируемая документация OpenAPI\,3.0 / Swagger UI на \code{/docs}.
  \item Все endpoints задокументированы: request/response schemas, примеры, коды ошибок.
\end{itemize}

\subsection{Инфраструктура}

\begin{itemize}[nosep]
  \item \textbf{Docker Compose} --- единственный способ запуска. \code{docker-compose up} поднимает все сервисы: simulator, backend, frontend, TimescaleDB, Redis.
  \item Конфигурация через \code{.env} файл с документированными переменными.
  \item \code{.env.example} в репозитории с описанием всех переменных.
  \item \code{README.md} с пошаговыми инструкциями запуска.
\end{itemize}

\subsection{Логи и наблюдаемость}

\begin{itemize}[nosep]
  \item Структурированные JSON-логи от всех backend-сервисов (\code{structlog}).
  \item Уровни логирования: DEBUG, INFO, WARNING, ERROR.
  \item Request logging с \code{trace\_id} для корреляции.
  \item Endpoint \code{/health} --- статус сервиса, связность БД, глубина очереди, uptime.
  \item Метрики: активные WebSocket-соединения, msg/sec, queue depth, HI computation time.
\end{itemize}


% =========================================================================
%  15. ПЛАН ТЕСТИРОВАНИЯ
% =========================================================================
\section{План тестирования}

\subsection{Unit Tests}

\begin{longtable}{l L{3cm} L{4.5cm} L{4.5cm}}
\toprule
\textbf{ID} & \textbf{Компонент} & \textbf{Тест} & \textbf{Ожидаемый результат} \\
\midrule
\endhead
UT-01 & HI Engine & Нормализация: значение в safe-зоне & Возвращает 1.0 \\
UT-02 & HI Engine & Нормализация: значение в warning-зоне & Пропорциональная оценка $(0, 1)$ \\
UT-03 & HI Engine & Нормализация: значение за critical & Возвращает 0.0 \\
UT-04 & HI Engine & Нормализация: значение NaN & Возвращает 0.0 \\
UT-05 & HI Engine & Штраф за алерты & Корректный мультипликативный штраф \\
UT-06 & HI Engine & Top-5 факторов: порядок & По убыванию contribution \\
UT-07 & HI Engine & Полный расчёт HI (профиль KZ8A) & Результат $\in [0, 100]$ \\
UT-08 & Ingest & Дедупликация duplicate timestamp & Второе сообщение отброшено \\
UT-09 & Ingest & Отклонение выброса & Значение вне физического диапазона отвергнуто \\
UT-10 & Ingest & EMA-сглаживание & Сглаженное значение соответствует $\alpha = 0.3$ \\
\bottomrule
\end{longtable}

\subsection{Integration Tests}

\begin{longtable}{l L{3.5cm} L{4cm} L{4cm}}
\toprule
\textbf{ID} & \textbf{Тест} & \textbf{Шаги} & \textbf{Ожидаемый результат} \\
\midrule
\endhead
IT-01 & Pipeline end-to-end & POST телеметрию $\to$ GET history & Данные в БД в течение 1\,с \\
IT-02 & WebSocket соединение & Подключиться, получить данные & JSON-фреймы с частотой 1\,Гц \\
IT-03 & WebSocket reconnect & Разорвать соединение, ждать & Клиент переподключается с backoff \\
IT-04 & Auth flow & Login с корректными creds & JWT-токен, доступ к protected endpoint \\
IT-05 & Auth rejection & Запрос без токена & Ответ \code{401 Unauthorized} \\
IT-06 & Role enforcement & Dispatcher $\to$ PUT config & Ответ \code{403 Forbidden} \\
IT-07 & Alert generation & Отправить значение выше порога & Алерт создан, отправлен через WS \\
IT-08 & Config update $\to$ HI & Изменить порог, проверить HI & HI пересчитан с новыми порогами \\
\bottomrule
\end{longtable}

\subsection{Scenario / E2E Tests}

\begin{longtable}{l L{3.5cm} C{1.5cm} L{5.5cm}}
\toprule
\textbf{ID} & \textbf{Сценарий} & \textbf{Время} & \textbf{Критерий успеха} \\
\midrule
\endhead
ST-01 & Normal operation & 5\,мин & Dashboard работает, HI стабильно 80--95, нет ложных алертов \\
ST-02 & Постепенный перегрев & 3\,мин & HI плавно снижается, алерты появляются, top factors показывают температуру \\
ST-03 & Внезапный сбой & 1\,мин & HI резко падает, статус Critical, множественные алерты, рекомендации \\
ST-04 & Highload burst $\times$10 & 1\,мин & UI не зависает, данные eventually consistent \\
ST-05 & Connection loss + recovery & 2\,мин & <<Нет связи>> через 3\,с, авто-reconnect, данные возобновляются \\
ST-06 & Replay mode & 2\,мин & Историческое воспроизведение с синхронизацией графиков/карты/алертов \\
ST-07 & Report export & 1\,мин & PDF/CSV скачивается с корректными данными \\
ST-08 & Переключение темы & 30\,с & Все экраны корректно отображаются в обеих темах \\
ST-09 & Смена типа локомотива & 1\,мин & Переключение KZ8A $\leftrightarrow$ ТЭ33А, параметры обновляются \\
\bottomrule
\end{longtable}

\subsection{Performance Tests}

\begin{tabularx}{\textwidth}{l L{5cm} l l}
\toprule
\textbf{ID} & \textbf{Тест} & \textbf{Целевое значение} & \textbf{Инструмент} \\
\midrule
PT-01 & WebSocket latency (end-to-end) & $< 500$\,мс & Сравнение timestamp \\
PT-02 & API response (history, 1\,ч) & $< 200$\,мс & cURL timing \\
PT-03 & Burst handling ($\times$10, 30\,с) & Без crash & Simulator burst mode \\
\bottomrule
\end{tabularx}


% =========================================================================
%  16. ФОРМАТ СДАЧИ И ДЕМОНСТРАЦИЯ
% =========================================================================
\section{Формат сдачи и демонстрация}

\subsection{Git-репозиторий}

\begin{itemize}[nosep]
  \item Публичный или общий Git-репозиторий.
  \item Чистая история (без секретов в коммитах).
  \item \code{README.md}: описание проекта, обзор архитектуры, prerequisites (Docker, docker-compose), пошаговая инструкция запуска (\code{docker-compose up}), учётные данные для демо-входа, скриншоты/GIF работающего приложения.
  \item \code{.env.example}: все необходимые переменные окружения с документацией.
\end{itemize}

\subsection{Презентация}

10--12 слайдов:
\begin{enumerate}[nosep]
  \item Постановка проблемы
  \item Архитектурная диаграмма
  \item Разбор формулы Health Index
  \item Обоснование выбора стека
  \item Демо-скриншоты ключевых экранов
  \item Вызовы и решения
\end{enumerate}

Формат: PDF или Google Slides.

\subsection{Демонстрация}

Live-демо или предзаписанное видео (3--5 минут), демонстрирующее:

\begin{itemize}[nosep]
  \item Нормальный режим работы с live-dashboard.
  \item Аварийный сценарий с деградацией HI и алертами.
  \item Highload burst ($\times$10).
  \item Потеря и восстановление соединения.
  \item Переключение темы.
  \item Replay mode.
  \item Смена типа локомотива.
\end{itemize}

\subsection{Архитектурная диаграмма}

Отдельный файл (PNG/SVG) в \code{docs/architecture.png}. Содержит все компоненты, потоки данных, порты и протоколы (см.~раздел~\ref{sec:arch-diagram}).


% =========================================================================
%  17. DEFINITION OF DONE
% =========================================================================
\section{Definition of Done}

\begin{enumerate}[nosep]
  \item Все сервисы поднимаются через единый \code{docker-compose up}.
  \item KZ8A и ТЭ33А поддерживаются через отдельные конфигурационные профили (без изменения кода).
  \item Swagger/OpenAPI документация доступна на \code{/docs}.
  \item Frontend отображает: dashboard с HI, панели узлов, карту маршрута, trends, replay.
  \item Светлая и тёмная тема полностью функциональны.
  \item Health Index считается в реальном времени с объяснимостью (top-5 факторов).
  \item Минимум 3 демо-сценария подготовлены (normal, fault, highload/connection loss).
  \item WebSocket reconnect с backoff + визуальная индикация <<нет связи>>.
  \item JWT-аутентификация с role-based access control.
  \item Endpoint \code{/health} возвращает статус сервиса.
  \item Структурированное JSON-логирование включено.
  \item Экспорт отчёта (PDF/CSV) функционален.
  \item Отметки событий (аннотации к алертам) реализованы.
  \item Презентация (10--12 слайдов) с архитектурой и формулой HI.
  \item \code{README.md} с полными инструкциями запуска.
\end{enumerate}


% =========================================================================
%  18. ГЛОССАРИЙ
% =========================================================================
\section{Глоссарий}

\begin{tabularx}{\textwidth}{l X}
\toprule
\textbf{Термин} & \textbf{Определение} \\
\midrule
HI & Health Index --- интегральный показатель состояния локомотива (0--100) \\
EMA & Exponential Moving Average --- экспоненциальное скользящее среднее \\
KZ8A & Электровоз казахстанского производства (переменный ток 25\,кВ) \\
ТЭ33А & Тепловоз с дизель-электрической тягой (GE Evolution) \\
NFR & Non-Functional Requirements --- нефункциональные требования \\
MVP & Minimum Viable Product --- минимально жизнеспособный продукт \\
SSE & Server-Sent Events --- односторонний серверный push через HTTP \\
WebSocket & Двунаправленный протокол связи поверх TCP \\
TimescaleDB & Расширение PostgreSQL для работы с time-series данными \\
Hypertable & Таблица TimescaleDB с автоматическим партиционированием по времени \\
JWT & JSON Web Token --- стандарт аутентификации на основе подписанных токенов \\
WCAG & Web Content Accessibility Guidelines --- стандарт доступности веб-контента \\
RBAC & Role-Based Access Control --- управление доступом на основе ролей \\
\bottomrule
\end{tabularx}


% =========================================================================
%  19. ПРИЛОЖЕНИЯ
% =========================================================================
\section{Приложения}

\subsection*{Приложение A. Параметры локомотивов}

Полная таблица телеметрических параметров для KZ8A и ТЭ33А (названия, единицы измерения, подсистемы, safe/warn/crit диапазоны, веса) --- \textit{см.\,\code{Locomotive\_Metrics\_Reference.md}}.

\subsection*{Приложение B. Формат JSON-сообщения телеметрии}

\begin{lstlisting}
{
  "loco_id": "KZ8A-001",
  "loco_type": "KZ8A",
  "ts": "2026-04-04T12:00:00.000Z",
  "parameters": {
    "speed_kmh": 87.3,
    "traction_motor_temp_c": 142.5,
    "voltage_v": 24800,
    "current_a": 350,
    "brake_pressure_bar": 4.8,
    "oil_temp_c": 78.2,
    "coolant_temp_c": 82.0,
    "fuel_level_pct": 65.0
  }
}
\end{lstlisting}

\subsection*{Приложение C. Формат WebSocket фрейма}

\begin{lstlisting}
{
  "type": "telemetry_update",
  "data": {
    "ts": "2026-04-04T12:00:01.000Z",
    "loco_id": "KZ8A-001",
    "parameters": { ... },
    "health_index": {
      "score": 82.5,
      "status": "Normal",
      "top_factors": [
        {
          "parameter": "traction_motor_temp_c",
          "subsystem": "Traction",
          "value": 142.5,
          "safe_range": [80, 130],
          "contribution": 0.15,
          "action": "Monitor temperature trend"
        }
      ],
      "subsystem_scores": {
        "Traction": 0.78,
        "Electrical": 0.95,
        "Cooling": 0.88,
        "Braking": 0.92
      }
    },
    "new_alerts": []
  }
}
\end{lstlisting}

\subsection*{Приложение D. Docker Compose --- обзор сервисов}

\begin{tabularx}{\textwidth}{l l l X}
\toprule
\textbf{Сервис} & \textbf{Образ} & \textbf{Порт} & \textbf{Зависимости} \\
\midrule
\code{frontend} & Node 20 + Next.js & 3000 & backend \\
\code{backend} & Python 3.11 + FastAPI & 8000 & db, redis \\
\code{simulator} & Python 3.11 & --- & redis \\
\code{db} & timescale/timescaledb & 5432 & --- \\
\code{redis} & redis:7-alpine & 6379 & --- \\
\bottomrule
\end{tabularx}

\subsection*{Приложение E. Цветовая палитра}

Конкретные цвета выбираются командой дизайна. Общие требования к палитре:

\begin{itemize}[nosep]
  \item Поддержка светлой и тёмной темы с достаточным контрастом (WCAG~AA).
  \item Семантические цвета статусов: зелёный (Normal), жёлтый/оранжевый (Attention), красный (Critical) --- должны быть интуитивно различимы.
  \item Единый акцентный цвет для интерактивных элементов.
  \item Иерархия текста: primary и secondary цвета для основного и вспомогательного текста.
  \item Нейтральные фоновые и поверхностные цвета для карточек, панелей и разделителей.
\end{itemize}


\end{document}
